\section{TDD：红绿重构}

\subsection{三阶段纪律}

\begin{importantbox}{红-绿-重构}
\begin{enumerate}
    \item \textbf{红（Red）}：先写一个\textbf{会失败}的测试——因为功能尚未实现，测试必然报错。
    \item \textbf{绿（Green）}：写\textbf{最少量}的代码使测试通过。
    \item \textbf{重构（Refactor）}：在保持绿色的前提下优化代码结构。
\end{enumerate}
\end{importantbox}

\subsection{为什么必须"亲眼看着测试失败"？}

\begin{warningbox}{事后测试证明不了任何事情}
\begin{itemize}
    \item 写完测试后必须运行，并看到\textbf{预期的报错}。
    \item 如果测试直接通过，说明测的东西是错的，或测的是已经存在的行为。
    \item 先写代码再补测试，往往只会测试"代码现在的行为"，而不是"代码应该的行为"——这就是\textbf{技术债务}的来源。
\end{itemize}
\end{warningbox}

\subsection{验证前必须有证据}

在宣称完成或测试通过前，必须\textbf{真实运行命令/工具}并检查输出——大量使用断言（assert），对结果有明确预期。

\subsection{本章小结}

TDD 的核心纪律是"先红后绿再重构"，强制 Agent 在每一步都用真实的测试结果作为证据，杜绝"假装完成"。


